%------------------------------------------------------------------------------%
\documentclass[a4paper,12pt,fleqn]{article}

\usepackage{calculo_varias_variaveis-1}

\ShowAnswers

%------------------------------------------------------------------------------%
\begin{document}

\makeHeader{CFVVI}{Exame -- Turma 7}{12/01/2026}

% 01
%------------------------------------------------------------------------------%
\exercise{40\%}
Determine os pontos candidatos a extremos de
\(
    x^2 + \ln y
\)
com $y>0$, sujeita à restrição
\(
    2x + y = 3
\)
\clearpagequestiononly

\begin{answer}
\begin{multicols}{2}
Queremos os extremos de
\[
    f(x, y) = x^2+\ln y \qquad y > 0
\]
sujeita à restrição
\[
    g(x, y) = 2x + y = 3
\]
Gradientes
\[
    \nabla f = \vetor{2x}{\sfrac{1}{y}}
    \qquad
    \nabla g = \vetor{2}{1}
\]

Pelo método dos multiplicadores de Lagrange
\begin{align*}
    2x          & = 2\lambda \\
    \frac{1}{y} & = \lambda \\
    2x + y      & = 3
\end{align*}
Da primeira equação temos que \(x=\lambda\), \\
substituindo na segunda temos
\[
    \frac{1}{y} = \lambda = x
    \qquad
    y = \frac{1}{x}
\]
Substituindo na restrição
\begin{align*}
    2x + y           & = 3  \\
    2x + \frac{1}{x} & = 3  \\
    2x^2 + 1         & = 3x \\
    2x^2 - 3x + 1    & = 0
\end{align*}
Resolvendo
\[
    \Delta
    = b^2 - 4ac
    = (-3)^2 - 4 \times 2 \times 1
    = 9 - 8
    = 1
\]
\[
    x
    = \frac{-b\pm\sqrt{\Delta}}{2a}
    = \frac{3\pm 1}{4}
\]

\smallskip
Para \(x=\frac{3+1}{4}=1\) temos \(y=1\)

\smallskip
Para \(x=\frac{3-1}{4}=\frac{1}{2}\) temos \(y=2\)

\smallskip
Os pontos candidatos a extremos são
\[
  P_1 = \left(1, 1\right)
  \quad \text{ e } \quad
  P_2 = \left(\frac{1}{2}, 2\right)
\]
\end{multicols}
\end{answer}

% 02
%------------------------------------------------------------------------------%
\begin{multicols}{2}
\raggedbottom
\exercise{20\%}
  Uma partícula descreve uma trajetória elíptica descrita
  pelas equações paramétricas
  \begin{gather*}
    x(t) = 3 \cos\left(\frac{\pi}{2}-t\right) + 2 \\
    y(t) = 5 \sin\left(\frac{\pi}{2}-t\right) + 1 \\
    0 \leq t \leq \pi
  \end{gather*}
  \begin{enumerate}[label=\alph*),itemsep=0pt,topsep=0pt]
    \item\!\! [20\%] calcule o vetor tangente à curva
    \item\!\! [10\%] avalie o vetor tangente em \\
    \hspace*{29pt} $t=0$, $t=\frac{\pi}{2}$ e em $t=\pi$
    \item\!\! [20\%] esboce a curva e os vetores \\
           \hspace*{29pt} tangentes calculados
  \end{enumerate}
\columnbreak
~

\begin{questiononly}
  \hfill\includegraphics{axis_xy}
\end{questiononly}
\begin{answeronly}
  \hfill\includegraphics{T2_A_parametrica}
\end{answeronly}
\end{multicols}
\clearpagequestiononly

\begin{answer}
\textbf{a)} \quad
O vetor tangente é
\(
  v(t) = \Vetor{\frac{dx}{dt}}{\frac{dy}{dt}}
\)

Avaliando cada derivada

\medskip
\[
  \frac{dx}{dt}
  = \frac{d}{dt}\left[2 + 3 \cos\left(\frac{\pi}{2}-t\right)\right]
  = 3 \frac{d}{dt}\cos\left(\frac{\pi}{2}-t\right)
  = -3 \sin\left(\frac{\pi}{2}-t\right)\frac{d}{dt}\left(\frac{\pi}{2}-t\right)
  = 3 \sin\left(\frac{\pi}{2}-t\right)
\]

\medskip
\[
  \frac{dy}{dt}
  = \frac{d}{dt}\left[1 + 5 \sin\left(\frac{\pi}{2}-t\right)\right]
  = 5 \frac{d}{dt}\sin\left(\frac{\pi}{2}-t\right)
  = 5 \cos\left(\frac{\pi}{2}-t\right)\frac{d}{dt}\left(\frac{\pi}{2}-t\right)
  = -5 \cos\left(\frac{\pi}{2}-t\right)
\]

\medskip
Portanto,
\(
  v(t) = \Vetor{3 \sin\left(\frac{\pi}{2}-t\right)}{-5 \cos\left(\frac{\pi}{2}-t\right)}
\)

\clearpage
\textbf{b)} \quad
Avaliando o vetor tangente nos pontos $t=0$, $t=\frac{\pi}{2}$ e em $t=\pi$
\setlength{\mathindent}{0pt}
\begin{multicols}{3}
\begin{align*}
  v(0)
  & = \Vetor{3 \sin\left(\frac{\pi}{2}-t\right)}{-5 \cos\left(\frac{\pi}{2}-t\right)} \evalat{0}{} \\[2mm]
  & = \Vetor{3 \sin\left(\frac{\pi}{2}\right)}{-5 \cos\left(\frac{\pi}{2}\right)} \\[2mm]
  & = \vetor{3 \times 1}{-5 \times 0} \\[2mm]
  & = \vetor{3}{0}
\end{align*}

\newcolumn
\begin{align*}
  v\left(\frac{\pi}{2}\right)
  & = \Vetor{3 \sin\left(\frac{\pi}{2}-t\right)}{-5 \cos\left(\frac{\pi}{2}-t\right)} \evalat{\frac{\pi}{2}}{} \\[2mm]
  & = \Vetor{3 \sin\left(\frac{\pi}{2}-\frac{\pi}{2}\right)}{-5 \cos\left(\frac{\pi}{2}-\frac{\pi}{2}\right)} \\[2mm]
  & = \Vetor{3 \sin\left(0\right)}{-5 \cos\left(0\right)} \\[2mm]
  & = \vetor{3 \times 0}{-5 \times 1} \\[2mm]
  & = \vetor{0}{-5}
\end{align*}

\newcolumn
\begin{align*}
  v(\pi)
  & = \Vetor{3 \sin\left(\frac{\pi}{2}-t\right)}{-5 \cos\left(\frac{\pi}{2}-t\right)} \evalat{\pi}{} \\[2mm]
  & = \Vetor{3 \sin\left(\frac{\pi}{2}-\pi\right)}{-5 \cos\left(\frac{\pi}{2}-\pi\right)} \\[2mm]
  & = \Vetor{3 \sin\left(\frac{-\pi}{2}\right)}{-5 \cos\left(-\frac{\pi}{2}\right)} \\[2mm]
  & = \vetor{3(-1)}{-5 \times 0} \\[2mm]
  & = \vetor{-3}{0}
\end{align*}
\end{multicols}

\bigskip
\textbf{c)} \quad
Calculando os pontos associados a $t=0$, $t=\frac{\pi}{2}$ e em $t=\pi$
\[
  x(0)
  = \left(3 \cos\left(\frac{\pi}{2}-t\right) + 2\right)\evalat{0}{}
  = 3 \cos\left(\frac{\pi}{2}\right) + 2
  = 3 \times 0 + 2
  = 2
\]
\[
  y(0)
  = \left(5 \sin\left(\frac{\pi}{2}-t\right) + 1\right)\evalat{0}{}
  = 5 \sin\left(\frac{\pi}{2}\right) + 1
  = 5 \times 1 + 1
  = 6
\]

\medskip
\[
  x\left(\frac{\pi}{2}\right)
  = \left(3 \cos\left(\frac{\pi}{2}-t\right) + 2\right)\evalat{\frac{\pi}{2}}{}
  = 3 \cos\left(0\right) + 2
  = 3 \times 1 + 2
  = 5
\]
\[
  y\left(\frac{\pi}{2}\right)
  = \left(5 \sin\left(\frac{\pi}{2}-t\right) + 1\right)\evalat{\frac{\pi}{2}}{}
  = 5 \sin\left(0\right) + 1
  = 5 \times 0 + 1
  = 1
\]

\medskip
\[
  x(\pi)
  = \left(3 \cos\left(\frac{\pi}{2}-t\right) + 2\right)\evalat{\pi}{}
  = 3 \cos\left(-\frac{\pi}{2}\right) + 2
  = 3 \times 0 + 2
  = 2
\]
\[
  y(\pi)
  = \left(5 \sin\left(\frac{\pi}{2}-t\right) + 1\right)\evalat{\pi}{}
  = 5 \sin\left(-\frac{\pi}{2}\right) + 1
  = 5 (-1) + 1
  = -4
\]
\end{answer}

% 03
%------------------------------------------------------------------------------%
\exercise{40\%}
Determine e classifique os pontos críticos da função
\(
  f(x,y) = e^{xy} - x^2 - y^2
\)

\clearpagequestiononly

\begin{answer}
  Calculando o gradiente
  \[
    f_x(x, y)
    = y e^{xy} - 2x
  \]
  \[
    f_y(x, y)
    = x e^{xy} - 2y
  \]
  Igualando o gradiente ao vetor nulo
  \[
    \begin{cases}
      y e^{xy} - 2x = 0\\
      x e^{xy} - 2y = 0
    \end{cases}
  \]
  rearranjando
  \[
    \begin{cases}
      2x = y e^{xy} \\
      2y = x e^{xy}
    \end{cases}
  \]
  Subtraindo as equações
  \[
    2(x-y) = -(x-y)e^{xy}
  \]

  Caso 1: $x\neq y$, dividimos os dois lados por $x-y$, obtendo
  \(
    2 = -e^{xy}
  \)
  que não possui solução.

  Caso 2: $x=y$, a equação \(2(x-y) = -(x-y)e^{xy}\) é satisfeita

  Substituindo na primeira equação e obtemos
  \begin{align*}
    2x & = y e^{xy} \\
    2x & = x e^{x^2}
  \end{align*}

  Se $x=0$, essa equação é satisfeita e temos o primeiro ponto crítico \(P_1=(0, 0)\)

  Se $x\neq y$, dividimos os dois termos por $x$
  \begin{align*}
    2x      & = x e^{x^2}        \\
    2       & = e^{x^2}          \\
    \ln(2)  & = x^2              \\
    \abs{x} & = \sqrt{\ln(2)}    \\
    x       & = \pm\sqrt{\ln(2)}
  \end{align*}
  Encontramos os pontos críticos
  \(P_2 = \left( \sqrt{\ln(2)},  \sqrt{\ln(2)}\right)\) e
  \(P_3 = \left(-\sqrt{\ln(2)}, -\sqrt{\ln(2)}\right)\)

  \bigskip
  Calculando a Hessiana de $f$
  \[
    f_{xx}(x, y)
    = y^2 e^{xy}-2
  \]
  \[
    f_{yy}(x, y)
    = x^2 e^{xy}-2
  \]
  \[
    f_{xy}(x, y)
    = xy e^{xy}+e^{xy}
  \]
  Discriminante (Determinante da Hessiana)
  \[
    D(x, y) = f_{xx}(x, y)f_{yy}(x, y) - f^2_{xy}(x, y)
  \]

  \bigskip
  Classificando o ponto \((0, 0)\)
  \[
    f_{xx}(0, 0) = -2
  \]
  \[
    f_{yy}(0, 0) = -2
  \]
  \[
    f_{xy}(0, 0) =  1
  \]
  \[
    D(0, 0) = 4-1 = 3 > 0
  \]
  Como \(f_{xx}(0, 0) = -2 < 0\), \textbf{o ponto \((0,0)\) é um máximo local}

  \bigskip
  No ponto \(\left(\sqrt{\ln 2}, \sqrt{\ln 2}\right)\)
  \[
    f_{xx}\left(\sqrt{\ln 2}, \sqrt{\ln 2}\right)
    = 2\ln 2 -2
  \]
  \[
    f_{yy}\left(\sqrt{\ln 2}, \sqrt{\ln 2}\right)
    = 2\ln 2 -2
  \]
  \[
    f_{xy}\left(\sqrt{\ln 2}, \sqrt{\ln 2}\right)
    = 2\ln 2 +2
  \]
  \[
    D\left(\sqrt{\ln 2}, \sqrt{\ln 2}\right)
    = (2\ln 2 -2)^2 - (2\ln 2 +2)^2
    < 0
  \]
  Portanto, é um \textbf{ponto de sela}

  \bigskip
  No ponto \(\left(-\sqrt{\ln 2}, -\sqrt{\ln 2}\right)\)
  \[
    f_{xx}\left(-\sqrt{\ln 2}, -\sqrt{\ln 2}\right)
    = 2\ln 2 -2
  \]
  \[
    f_{yy}\left(-\sqrt{\ln 2}, -\sqrt{\ln 2}\right)
    = 2\ln 2 -2
  \]
  \[
    f_{xy}\left(-\sqrt{\ln 2}, -\sqrt{\ln 2}\right)
    = 2\ln 2 +2
  \]
  \[
    D\left(-\sqrt{\ln 2}, -\sqrt{\ln 2}\right)
    = (2\ln 2 -2)^2 - (2\ln 2 +2)^2
    < 0
  \]
  Portanto, é um \textbf{ponto de sela}
\end{answer}

%------------------------------------------------------------------------------%
\end{document}
%------------------------------------------------------------------------------%


